\section{Soot formation at high fuel loadings}
\label{sec:paper2_high_fuel_loadings}


This section presents the paper titled ``\textit{Gas-phase chemistry and soot formation during shock-tube pyrolysis of methane at high fuel loadings}'' and accepted in the 41st International Symposium on Combustion, where Omnisoot was utilized to extend the application of kinetic mechanisms from dilute to high-fuel-loading regime typical in cogeneration of carbon black and hydrogen. The paper provides the first simultaneous investigation of methane pyrolysis chemistry and particulate formation under high fuel concentrations accompanied with numerical simulations that accounts for the evolution of gas chemistry and formation of particulate. Detailed measurement-model comparisons enabled the assessment of several kinetic mechanisms to elucidate the combined effect of temperature and fuel loading on methane carbon flux to intermediates that determine the yield of carbonaceous particles.

Direct decomposition of methane offers a low-emission route for large-scale co-production of carbon black (CB) and hydrogen~\citep{fulcheri2024plasma}. Achieving CB with desired specific surface area, morphology, and internal structure requires precise control of gas-phase chemistry as well as CB inception and surface growth~\citep{kang2025effect}. This remains challenging due to the complex reaction pathways~\citep{johansson2018resonance} and short time scales~\citep{martin2022soot} involved in methane-to-CB conversion. Although kinetic mechanisms coupled with particle dynamics have been developed to describe polycyclic aromatic hydrocarbon (PAH) formation and particle inception~\citep{frenklach2020mechanism}, most were developed for dilute conditions to facilitate the inference of rate constants~\citep{wang2016improved}, and to minimize spectral interference in laser diagnostics~\citep{kohse2005combustion}. However, high fuel loadings (local mole fractions $\sim$1) are typical in industrial CB synthesis~\citep{fulcheri2023energy}, highlighting the need for data to benchmark kinetic mechanisms across low- to high-fuel regimes. 

Prior studies on soot formation from methane at high fuel loadings in flames~\citep{liu2022effect} and flow reactors~\citep{mokashi2024understanding} lack time- or space-resolved species data due to diagnostic challenges such as strong radiation, interference from intermediates, and limited optical access~\citep{karatacs2012soot}. Time-resolved CH$_4$ measurements during 10\%~CH$_4$ pyrolysis were initially conducted to refine decomposition reaction rates~\citep{nativel2019shock}. Later, simultaneous measurements of $\mathrm{CH_4}$, $\mathrm{C_2H_4}$, and $\mathrm{C_2H_2}$ were performed for 10–40\% $\mathrm{CH_4}$ pyrolysis with a focus on maximizing H$_2$ yield without soot characterization~\cite{ferris2024experimental}. 

In this work, we investigate the shock-tube pyrolysis of 5\%, 10\%, and 30\% $\mathrm{CH_4}$-Ar mixtures at  P$_5$=4$\pm$0.9~atm and 1800$<$T$_5$$<$2500~K, providing simultaneous measurements of key intermediates and soot volume fraction via optical diagnostics. Simulations were conducted using detailed kinetic mechanisms coupled with a sectional population balance model that accounts for inception, surface growth, and agglomeration. Comparison between measurements and model predictions enables (i) evaluation of kinetic mechanism performance for $\mathrm{CH_4}$ conversion and the resulting carbon flux governing $\mathrm{C_2H_4}$ and $\mathrm{C_2H_2}$ formation and (ii) assessment of soot volume fraction evolution over extended times to examine their dependence on temperature and fuel loading.
